\section{Discussion}\label{sec:discussion}

\subsection{What the result does and does not say}

The model does not imply that real-world flood protection, sea walls, drainage systems, or other adaptation investments are generally excessive. Such projects protect households, lives, public infrastructure, nonindustrial assets, and sectors outside the model. The analysis isolates an industrial-organization margin that standard project appraisal can miss: firms may reduce private geographic backup capacity in response to stronger place-based public protection, and decentralized governments may still value the same protection as a tool for attracting production.

The appropriate policy interpretation is therefore joint evaluation. A regional protection project that materially changes firms' location and continuity decisions should be evaluated together with those induced private organizational responses. The theorem identifies circumstances in which the marginal industrial-resilience component of a project can become negative even though the engineering reliability effect remains positive.

\subsection{Why the model is kept minimal}

Several extensions are natural but intentionally excluded from the baseline. Firms could choose multiple backup sites, governments could offer tax subsidies, disasters could be spatially correlated across jurisdictions, or the product market could use price rather than quantity competition. These may matter for quantitative applications, but none is required to state the present mechanism. Adding them before establishing the benchmark would obscure the distinction between the paper's contribution and established results on fiscal competition, adaptation crowd-out, or resilient supply networks.

The reduced-form backup choice should likewise not be read as a literal probability purchased from a technology. It is a tractable index of costly preparedness that increases the probability of maintaining production after primary-site disruption. The institutional examples in Section~\ref{sec:institutions} show that firms have multiple practical margins that fit this interpretation.

\subsection{Generality}

The same strategic logic can apply outside manufacturing whenever four ingredients are present: location-specific correlated disruption risk, public investment that lowers local risk, mobile competing firms, and costly private geographic redundancy. Data centers with geographic failover capacity and logistics operators with alternative terminal or routing capacity are possible examples. These interpretations require no new strategic force, but the paper keeps manufacturing as its canonical application.
